\chapter{A computational method in Plasma physics}

\section{The Variational Principle}

\subsection{The Magnetostatic Equation}

Variational priciple or the potential energy

\[
E = \int \int \int \left(\frac{B^2}{2} + \frac{p}{\gamma -1 }   \right)
  \]

  In minimizing the potential energy we require that $\grad \cdot B = 0$ everywhere, the toroidal and poloidal fluxes within each surface are fixed, mass within each flux tube has a fixed value, the total toroidal and poloidal fluxes in the vacuum are fixed, and the free surface $\Gamma$ and outer wall are flux surfaces wich $B \cdot v = 0$


\subsection{Flux Constraints in the Plasma}

Taking
\[
B = \grad s \times \grad \psi
  \]
  and knowing that B must be single valued, we take the special form
\[
\psi = -u+ \mu(s) v + \lambda(s,u,v)
  \]
  Though a more general form is often used

  \[
\psi = f_1(s)u+ f_2(s) v + \lambda(s,u,v)
  \]
  where $\lambda$ is periodic in $u$ and $v$. We then find the toroidal flux
\[
F_T(s_0) = \oint \psi ds
  \]
\subsection{The Ergodic Constraint}

Taking the relationship that
\[
B \cdot \grad p = 0
  \]

  Implies that the magnetic lines are real characteristics of the magnetostatic equation. In order to formulate a well-posed problem, we want to eliminate the real characteristics from our system of equations.
